\section{Ograniczenia}
\subsection{Ogólne}
Przepływ musi być nieujemny oraz mniejszy przepustowości
\begin{equation}
0 \leq f_{ij} \leq u_{ij}, \qquad (i,j) \in N \\
\end{equation}
Wszystko co przypłyneło do stanu k musi stamtąd "odpłynąć"
\begin{equation}
\sum_{i \in N \backslash \{z\}} f_{ki} = \sum_{i \in N \backslash \{W\}} f_{jk}, \qquad k \in N \backslash \{z, W\}
\end{equation}

\subsection{Dostępne surowce}
Nie możemy kupić więcej surowca S1 niż jest dostępne
\begin{equation}
f_{zS1} \leq 12000
\end{equation}
Nie możemy kupić więcej surowca S2 niż jest dostępne
\begin{equation}
f_{zS2} \leq 8000
\end{equation}

\subsection{Koszt zakupu surowców}
Nie możemy kupić mniej niż 0 surowców
\begin{equation}
a_1 \geq 0
\end{equation}
\begin{equation}
a_2 \geq 0
\end{equation}
\begin{equation}
b_1 \geq 0
\end{equation}
\begin{equation}
b_2 \geq 0
\end{equation}
Ograniczamy ile ton ponad dany limit kupiliśmy
\begin{equation}
a_1 \leq f_{zS1} - 2387
\end{equation}
liczba powyżej 2056 ton
\begin{equation}
a_2 \leq f_{zS1} - 6659
\end{equation}
liczba powyżej 2056 ton
\begin{equation}
b_1 \leq f_{zS2} - 2090
\end{equation}
liczba powyżej 2056 ton
\begin{equation}
b_2 \leq f_{zS2} - 4349
\end{equation}
\subsection{Transport surowca S1 do przygotowalni}
wagon może przetransportować maks. 18 ton surowca S1
\begin{equation}
b_{P1} \geq [\frac{ f_{S1, G1} }{18}]
\end{equation}
\subsection{Transport surowca S2 do przygtowalni}
ciężarówka może przetransportować maks. 25 ton surowca S2
\begin{equation}
c_{P2} \geq [ \frac{ f_{S2, P2} }{ 25 } ]
\end{equation}
\subsection{Praca przygotowalni}
Przygotowalnia nie może mieć więcej surowców niż przepustowości
\begin{equation}
\sum_{i \in N, p \in P} f_{ip} \leq 16000
\end{equation}
ilość półproduktów w zależności od surowca
\begin{equation}
(\sum_{i \in N} f_{ip}) m_{pd} = f_{pd}, \qquad i \in N, p \in P, d \in D
\end{equation}
\subsection{Koszt pracy przygotowalni}
3 pracowników może obsługiwać maksymalnie 150 ton całkowitej ilości przerabianych surowców
\begin{equation}
l \geq [ \frac{ \sum_{i\in N, p \in P} f_{ip} }{150} ]
\end{equation}
\subsection{Transport S2 do zakładu obróbki cieplnej}
ciężarówka może przetransportować maks. 25 ton surowca S2
\begin{equation}
c_{O2} \geq [\frac{ f_{S2, O2} }{25}]
\end{equation}
\subsection{Praca zakładu obróbki cieplnej}
ilość surowca S2 w zakładzie obróbki cieplnej nie może być większa od 6000
\begin{equation}
f_{S2, O2} \leq 6000
\end{equation}
\subsection{Minimalna dostarczona ilość wyrobów}
Zgodnie z umową musimy dostarczyć przynajmniej 5000 ton wyrobów
\begin{equation}
\sum_{i \in N f_{iw}} \geq 5000, \qquad w \in W
\end{equation}
\subsection{Maksymalna liczba wagonow}
Na jedna lokomotywę, nie może przypadać więcej niż 12 wagonów
\begin{equation}
12wl \geq c_{P1}
\end{equation}

